\SourceLeaf{355}

\begin{center}
\large\bfseries GEOMETRISCHE SEITE DER TERNÄREN FORMEN.
\end{center}



Ein Punkt im Raume $(0)$ sei als Anfangspunkt angenommen. Der Übergang von da zu drei
andern Punkten $\boldsymbol{P}$, $\boldsymbol{P}'$, $\boldsymbol{P}''$, die mit jenem nicht
in einer Ebene liegen, sei resp. $t$, $t'$, $t''$; wo, so oft keine Verwechslung möglich
ist, die Punkte $\boldsymbol{P}$, $\boldsymbol{P}'$, $\boldsymbol{P}''$ selbst durch $(t)$,
$(t')$, $(t'')$ bezeichnet werden mögen.

Es sei ferner allgemein $(t,t')$ das Product der Länge (kürzesten Entfernung) der beiden
Linien $t$, $t'$ in den Cosinus ihrer Neigung etc.

Man hat allgemein $(\alpha t + \alpha' t' + \alpha'' t'' + .\,.\,.,\ \beta u + \beta' u' +
\beta'' u'' + .\,.\,.)$, wenn man die Multiplication
\[
(\alpha t + \alpha' t' + \alpha'' t'' + .\,.\,.) \times (\beta u + \beta' u' + \beta'' u'' + .\,.\,.)
\]
ausführt und statt $tu$, $tu'$, $tu''$, $t'u'$, $t'u''$ u.\,s.\,w. $(t,u)$, $(t,u')$,
$(t,u'')$, $(t',u')$, $(t',u'')$ u.\,s.\,w. schreibt.

Jeder Punkt im Raume wird durch ein Trinomium
\[
(xt + x't' + x''t'')
\]
dargestellt werden können.

Für alle Punkte, die in einer bestimmten Ebene liegen, wird dann eine Gleichung
\[
\lambda x + \lambda' x' + \lambda'' x'' = L
\]
stattfinden, wo $\lambda$, $\lambda'$, $\lambda''$, $L$ bestimmte Zahlen bedeuten. Für eine
Ebene durch die drei Punkte $\mu t$, $\mu' t'$, $\mu'' t''$ ist
\[
\lambda\mu = \lambda'\mu' = \lambda''\mu'' = L.
\]
